% ルビ（振り仮名）つきの和文。
%
% ルビは本文の**すぐ上に重ねて**組まれるので、行にまとめる段で本文の行に紛れ込む。
% 直す前の実測:
%
%   `そ かく かたまり本論文では疎な注意機構の核について論じる。`
%     ← 1 行ぶんのルビが全部**文の先頭に固まる**（ルビが別の行になった場合）
%   `そ本手法の中心となる考え方は、…`
%     ← 別の行のルビが**次の文の頭**に付く（ルビが前の行に吸われた場合）
%
% どちらも原文に無い文字列が本文に入る。ルビは 7 文字あることもあるので
% 「置き場所を間違えた断片」（4 文字以下）の網にも掛からず、印も申告も出ない。
%
% 正解には漢字の側だけを残す（`\ruby{漢字}{よみ}` の 1 つ目）。読み仮名が
% 本文に混ざったら壊れている。
%
% 正解の取り出し方は他の real-papers/*.tex と同じ（本文 1 文 = \s{...} を 1 行）。
\documentclass[11pt,a4paper]{article}

\usepackage{fontspec}
\setmainfont{Noto Serif CJK JP}
\XeTeXlinebreaklocale "ja"
\XeTeXlinebreakskip=0pt plus 0.4pt minus 0.4pt
\usepackage[margin=2.4cm]{geometry}

% 本文の上に小さく重ねるルビ。実物の \ruby と同じ見え方にする。
\newcommand{\ruby}[2]{\leavevmode\hbox{\ooalign{\hfil #1\hfil\crcr
  \raise 1.15em\hbox to 0pt{\hss\fontsize{5.5}{5.5}\selectfont #2\hss}}}}

\newcommand{\s}[1]{#1}

\title{注意機構の疎な核について}
\author{田中 芽衣}
\date{}

\begin{document}
\maketitle

\section{はじめに}

\s{本論文では\ruby{疎}{そ}な注意機構の\ruby{核}{かく}について論じる。}
\s{実験では固定の記憶容量の下で、\ruby{塊}{かたまり}の大きさと到達可能な処理速度の間にどのような釣り合いがあるかを測った。}
\s{提案手法は既存の手法に比べて精度を保ったまま遅延を三割ほど削減できることを示す。}
\s{以下では、まず前提を述べ、続いて評価の手順を説明し、最後に限界について議論する。}

\section{提案手法}

\s{本手法の中心となる考え方は、注意の重みが実際には\ruby{疎}{そ}であるという観察に基づいている。}
\s{これは経験的に確かめられており、多くの実装で同じ傾向が報告されている。}
\s{したがって、密な計算をすべて行う必要はない。}
\s{\ruby{塊}{かたまり}ごとに代表値を求めてから、上位のものだけを選ぶ手順を取る。}
\s{この手順は既存の実装に小さな変更を加えるだけで組み込める。}

\end{document}
